\documentclass[11pt]{article}

% Shared preamble for the SKJS soundness corpus (main paper + deep dives).
% Loaded via \input in each document's preamble. Documents may define
% additional packages, theorem environments, and macros AFTER this input,
% but must not redefine anything here.
\usepackage{amsmath}
\usepackage{amssymb}
\usepackage{amsthm}
\usepackage{mathtools}
\usepackage{stmaryrd}
\usepackage[margin=1in]{geometry}
\usepackage{tcolorbox}
\usepackage{booktabs}
\usepackage{array}
\usepackage{longtable}
\usepackage[hyphens]{url}
\urlstyle{tt}

% Long typewriter tokens (paths, identifiers) defeat hyphenation;
% allow inter-word stretch instead of overfull lines.
\emergencystretch=2em

\newtheorem{theorem}{Theorem}
\newtheorem{lemma}[theorem]{Lemma}
\newtheorem{definition}[theorem]{Definition}
\newtheorem{corollary}[theorem]{Corollary}
\newtheorem{axiom}[theorem]{Axiom}
\newtheorem{obligation}[theorem]{Obligation}

% Summary box for key definitions
\newtcolorbox{keydef}{
  colback=gray!5,
  colframe=gray!50,
  boxrule=0.5pt,
  arc=2pt,
  left=4pt,right=4pt,top=2pt,bottom=2pt
}

% Core semantic notation shared by every document
\newcommand{\sem}[1]{\llbracket #1 \rrbracket}
\newcommand{\Val}{\mathsf{Val}}

% Implementation anchor: file / class / function at the pinned revision
\newcommand{\impl}[1]{\par\noindent{\footnotesize\textsf{impl:} \texttt{#1}}}

% Single source for the modeled revision. Every document's \date uses this;
% the ledger and tooling read the same values.
\newcommand{\skjspinhash}{9108f339}
\newcommand{\skjspindate}{2026-07-02}
\newcommand{\skjspin}{Models revision \texttt{\skjspinhash} of \texttt{main} (\skjspindate)}


\usepackage{mathpartir}

\newtheorem{finding}[theorem]{Finding}

% Macros copied from the main paper (skjs-type-soundness.tex)
\newcommand{\SType}{\mathsf{SType}}
\newcommand{\STypes}{\mathsf{STypes}}
\newcommand{\TType}{\mathsf{TType}}
\newcommand{\TInst}{\mathsf{TInst}}
\newcommand{\TInsts}{\mathsf{TInsts}}
\newcommand{\TParam}{\mathsf{TParam}}
\newcommand{\Clauses}{\mathsf{Clauses}}
\newcommand{\Clause}{\mathsf{Clause}}
\newcommand{\Subst}{\mathsf{Subst}}
\newcommand{\SVOpen}{\mathsf{SVOpen}}
\newcommand{\SVClosed}{\mathsf{SVClosed}}
\newcommand{\clause}[2]{\langle #1 \mid #2 \rangle}

% Local macros
\newcommand{\expandop}{\mathsf{expand}}
\newcommand{\instop}{\mathsf{instantiate}}
\newcommand{\normop}{\mathsf{normalize}}
\newcommand{\cstr}{\mathsf{cstr}}
\newcommand{\Cov}{\mathsf{cov}}
\newcommand{\Contra}{\mathsf{contra}}
\newcommand{\NoneM}{\mathsf{none}}
\newcommand{\getop}{\mathsf{get}}
\newcommand{\refop}{\mathsf{refine}}

\title{Type Expansion in SKJS:\\
the Instance Layer, the Expansion Cache, and Axiom \textbf{L-EXP}\\
{\large Deep Dive C of the SKJS Soundness Series}}
\author{}
\date{\skjspin.\\
Companion to the main soundness paper.}

\begin{document}

\maketitle

\begin{abstract}
The main soundness paper assumes axiom \textbf{L-EXP}: expansion of a type
to its instance-plus-clauses view preserves or enlarges its semantics, with
residual instances keeping theirs. This deep dive models expansion precisely
as a worklist over the instance layer, proves the containment
$\sem{t}_\rho \subseteq \sem{\expandop(t)}_{\rho'}$ for the expansion modes
the refinement machinery actually uses, and — as its centerpiece — turns the
\texttt{seenExpansion} cache-soundness comment
(\texttt{TTypeSubst.sk:17--31}) into a precise invariance theorem. The
audit yields a result \emph{stronger} than the comment claims for type
aliases: normalization is \emph{binding-blind} — every substitution it
consults for variable bindings is freshly created empty — so a cached alias
expansion is a function of the instance alone and is immune to the
freeze-window staleness that threatens the other memo caches. It also
yields corrections: the comment's ``\texttt{inter} never resolves an
unbound \texttt{TVar}'' is false for variables nested in compound operands
(harmless where it matters, for a stated reason); the class-instantiation
path is \emph{not} binding-blind and remains an open obligation; and
contravariant-mode expansion does not satisfy \textbf{L-EXP} at all, forcing
a scope edit to the axiom. Section~\ref{sec:status} gives the resulting
ledger edits for the main paper's Assumptions table.
\end{abstract}

\tableofcontents

%==========================================================================
\section{Purpose}
\label{sec:purpose}

\paragraph{Row addressed.}
\textbf{L-EXP} (the main paper's Assumptions and
Proof Obligations table records the mode-scoped \textsc{proved} status
this paper establishes): ``Expansion preserves or enlarges semantics;
residual instances keep theirs,'' anchored at \texttt{TType::expand}. Every
refinement operator of the main paper runs on the expanded view
(its Definition ``Memberwise Refinement''), and the proof of its
memberwise-refinement lemma consumes \textbf{L-EXP} in its first step; the
axiom is therefore load-bearing for every branch-soundness result.

\paragraph{What this paper does.}
(i) models $\expandop$ as a function of a type, a substitution state, and a
variance mode, exactly following \texttt{TType::expandToVectors};
(ii) proves the containment per instance construct, splitting off two named
interface assumptions where the argument crosses into \texttt{TEnv}
(normalization) and array-literal commitment;
(iii) proves the \texttt{seenExpansion} cache-invariance theorem for the
type-alias path and names the class path as an obligation;
(iv) records the state evolutions expansion can cause and defers their
transport to the main paper's \textbf{L-MONO};
(v) states precise ledger edits (Section~\ref{sec:status}).
Failed proof attempts are reported as numbered Findings.

\paragraph{Conventions.}
Notation, semantics, and the status vocabulary
(\textsc{proved} / \textsc{assumed} / \textsc{open} / \textsc{violated})
are the main paper's. All code references are against the pinned revision
\texttt{\skjspinhash}; line numbers are cited where they disambiguate.
``Paper B'' is the deep dive on union and intersection
(\textbf{L-UNION}/\textbf{L-INTER}), which this paper cites and does not
re-prove.

%==========================================================================
\section{Minimal Recap}
\label{sec:recap}

\subsection{Representation and semantics}

A $\TType$ is a pair $(I, C)$ of an instance set $I \in \TInsts$ and
clauses $C \in \Clauses$, with
$\sem{(I,C)}_\rho = \sem{I}_\rho \cup \sem{C}_\rho$
(main paper, Definition ``Clause, Clauses, TType''). A clause
$\clause{V}{S}$ pairs type parameters with an $\STypes$ union or $\top$;
the top clause $\clause{\emptyset}{\top}$, if present, is alone. The
instance layer $\TInst$ has four constructors:
type variables $\alpha$ (\texttt{TVar}, carrying a finite set of
\emph{occurrence-level} constraint types — \texttt{TVar::setConstraint} is
a functional update returning a new \texttt{TVar} value with the same
identifier — registry row SETCONSTRAINT-FUN, deep dive E), unexpanded
named instantiations
\texttt{TSInst}$(n, \bar{a})$, unexpanded intersections
\texttt{TInter}$(\{t_1,\dots,t_k\})$, and delayed array literals
\texttt{TArrayLit}$(\alpha, \mathit{elts}, \mathit{fresh})$.
\impl{TType.sk / TType, TInsts, Clause, Clauses}
\impl{TType.sk / TVar (2576), TSInst (2626), TArrayLit (2618);
TTypeInter.sk / TInter (17)}

The main paper gives the instance layer an abstract, state-indexed
semantics. We spell out the part this paper consumes:

\begin{definition}[Instance Semantics Interface]
\label{def:inst-sem}
For a substitution state $\rho$ (Definition~\ref{def:subst-recap}):
\begin{itemize}
\item[(V1)] $\sem{\alpha^K}_\rho \subseteq \bigcap_{c \in K}\sem{c}_\rho$
  for a variable occurrence $\alpha^K$ with constraint set $K$;
\item[(V2)] if $\rho$ binds $\alpha$ (as $\SVOpen(t)$ or $\SVClosed(t)$),
  then $\sem{\alpha^K}_\rho \subseteq \sem{t}_\rho$;
\item[(I1)] $\sem{\mathsf{TInter}(\{t_i\})}_\rho =
  \bigcap_i \sem{t_i}_\rho$ (an unexpanded intersection \emph{means} the
  intersection);
\item[(N1)] $\sem{\mathsf{TSInst}(n,\bar{a})}_\rho$ is the denotation of
  the named definition $n$ at arguments $\bar{a}$ — related to expansion
  only through assumption \textbf{E-NORM} below;
\item[(A1)] $\sem{\mathsf{TArrayLit}(\alpha,\mathit{elts},\_)}_\rho$ is
  the set of runtime values of the originating array-literal expression —
  related to expansion only through assumption \textbf{E-ARR} below.
\end{itemize}
(V1)/(V2) are the main paper's canonical variable denotation, registry row
TVAR-SEM: a bound occurrence denotes its envelope intersected with its
occurrence constraints, and an unbound occurrence denotes $\emptyset$.
This paper imports that definition.
\end{definition}

\subsection{The substitution state, condensed}
\label{sec:subst-recap}

\begin{definition}[State; get transitions used by expansion]
\label{def:subst-recap}
$\rho$ maps \texttt{TVarID}s to $\SVOpen(t)$ or $\SVClosed(t)$ and carries
the flag $\mathsf{freeze}$ (\texttt{freezeCandidates}, \emph{not} part of
the saved state). Expansion reads variables through
$\getop(\alpha, \mathit{var})$:
covariant get on an \emph{unbound} $\alpha$ seals it to top
($\alpha \mapsto \SVClosed(\top)$, returning $\top$); on $\SVOpen(t)$ it
seals to $\SVClosed(t)$ and returns $t$ — except under $\mathsf{freeze}$
for a candidate, where it returns $t$ \emph{without} sealing; on
$\SVClosed(t)$ it returns $t$. Contravariant get returns $\bot$ on
unbound (no write) and the bound value otherwise. Both then intersect the
result with the occurrence's constraint types via \texttt{TType::inter}.
The full transition table is the main paper's Definition
``Transitions''.
\impl{TTypeSubst.sk / Subst::get (188), freezeCandidates (89)}
\end{definition}

The mutable \texttt{Subst} also carries four memo caches —
\texttt{seenSub}, \texttt{seenInter}, \texttt{seenUnion}, and this paper's
subject \texttt{seenExpansion} $:\ \mathsf{TSInst} \to \TType$ — all four
part of the saved state: \texttt{save}/\texttt{restore} and \texttt{clone}
copy bindings and caches \emph{together}.
\impl{TTypeSubst.sk / SubstState (13--49), SavedSubst (51), clone (137)}

\subsection{The axiom under study}

\begin{axiom}[L-EXP, as stated in the main paper]
\label{ax:lexp}
Expansion rewrites a type to an instance-plus-clauses view using the
current substitution state, possibly evolving it ($\rho \to \rho'$, e.g.\
sealing a fresh covariantly-read variable to top). Then
$\sem{t}_\rho \subseteq \sem{\expandop(t)}_{\rho'}$, and residual
instances keep their semantics.
\end{axiom}

Two things the main paper leaves open, which this paper fixes: expansion
takes a \emph{mode} argument (\texttt{expandTVars :\ ?Variance}), and the
containment holds for modes $\NoneM$ and $\Cov$ but \emph{not} for
$\Contra$ (Section~\ref{sec:tvar}); and ``residual instances'' means
exactly: unresolved type variables in mode $\NoneM$, plus unexpanded
recursion cut-offs re-emitted by \texttt{instantiate} and \texttt{TInter}
cycle guards (Sections~\ref{sec:tsinst}, \ref{sec:tinter},
\ref{sec:cycles}).

%==========================================================================
\section{Expansion as a Function}
\label{sec:model}

\subsection{The worklist}

\begin{definition}[Expansion]
\label{def:expand}
For a mode $m \in \{\NoneM, \Cov, \Contra\}$ (the implementation's
\texttt{expandTVars} $\in$
\{\texttt{None()}, \texttt{Some(Covariant())},
\texttt{Some(Contravariant())}\}), define
$\expandop_m(t, \rho) = (I', C', \rho')$ by the following worklist,
mirroring \texttt{TType::expandToVectors} line by line. If
$t$ is top, the result is $\top$ (returned as \texttt{None}). Otherwise
initialize $C' := t.\mathsf{clauses}$ (as a vector), $I' := \emptyset$,
$\mathit{todo} := t.\mathsf{insts}$, $\mathit{seen} := t.\mathsf{insts}$.
While $\mathit{todo}$ is nonempty, pop $i$ and compute a one-step
expansion $E_m(i)$:
\begin{itemize}
\item $i = \mathsf{TSInst}(n,\bar{a})$: if cached, $E_m(i) :=
  \rho.\texttt{seenExpansion}(i)$; else
  $E_m(i) := \instop(n, \bar{a})$ and the cache is extended
  ($\rho.\texttt{seenExpansion}[i] := E_m(i)$). Mode-independent.
\item $i = \mathsf{TInter}(\{t_j\})$: $E_m(i) :=
  \mathsf{interN}(\{t_j\})$ with \texttt{tenv.expandInter} set
  (Section~\ref{sec:tinter}). Mode-independent.
\item $i = \alpha^K$: in mode $\NoneM$, push $\alpha^K$ onto $I'$
  (\emph{residual}) and continue; in modes $\Cov$/$\Contra$,
  $E_m(i) := \getop(\alpha^K, m)$ — the binding read intersected with the
  occurrence constraints, with the seal side effects of
  Definition~\ref{def:subst-recap}.
\item $i = \mathsf{TArrayLit}(\alpha, \mathit{elts}, \_)$: $E_m(i)$ is
  the committed type $\mathsf{peek}(\alpha)$ if $\alpha$ is bound
  (no state change), else the mutable tuple type over $\mathit{elts}$
  (Section~\ref{sec:arraylit}). Mode-independent.
\end{itemize}
If $E_m(i)$ is top, the whole expansion returns $\top$. Otherwise extend
$C'$ with $E_m(i).\mathsf{clauses}$ and push each instance of
$E_m(i).\mathsf{insts}$ not in $\mathit{seen}$ onto $\mathit{todo}$
(adding it to $\mathit{seen}$). On termination,
$\expandop_m(t,\rho) = (I', C', \rho')$ where $\rho'$ is the evolved
state. \texttt{TType::expand} rebuilds the pair from
$\TInsts{::}\mathsf{unionN}(I')$ and
$\Clauses{::}\mathsf{unionN}(C')$;
\texttt{expandToClauses} does the same in mode
\texttt{Some(variance)} and returns the clause part.
\impl{TType.sk / expandToVectors (634), expand (680), expandToClauses
(692)}
\end{definition}

\begin{lemma}[State Evolutions of Expansion]
\label{lem:evolutions}
The only substitution evolutions $\expandop_m$ performs are:
(i) $\getop$ transitions — seal-to-top of unbound variables and
value-preserving seals $\SVOpen(t) \to \SVClosed(t)$ (modes
$\Cov$; mode $\Contra$ performs no binding writes); (ii) memo-cache
extensions (\texttt{seenExpansion} here; \texttt{seenInter} /
\texttt{seenUnion} inside the invoked operators); (iii) version bumps.
In particular expansion never executes a $\mathsf{grow}$: no binding
value is ever unioned or replaced.
\end{lemma}
\begin{proof}
By inspection of Definition~\ref{def:expand} against the code: the
\texttt{TVar} arm calls \texttt{Subst::get} only; the \texttt{TSInst}
arm calls \texttt{tenv.instantiate(subst, subst, ...)}, and
\texttt{instantiate} immediately re-packages the state into a fresh
\texttt{Subst::mcreate(tcenvb, baseSubst)} \emph{copy} — its binding
writes (constraint-check clones, class-body expansion) land in the copy
and are dropped; only the caller's \texttt{seenExpansion} write at
\texttt{TType.sk:653} touches the live state. The \texttt{TInter} and
\texttt{TArrayLit} arms call \texttt{interN} (get-transitions inside
\texttt{inter\_}'s covariant expansions, plus \texttt{seenInter}) and
\texttt{peek} (read-only). \qedhere
\impl{TEnv.sk / instantiate (110, 135); TType.sk / expandToVectors (631);
TTypeSubst.sk / peek (224)}
\end{proof}

Transport of already-installed types across these evolutions is exactly
the main paper's axiom \textbf{L-MONO} (whose anchor names the
\texttt{Subst::get} seal-to-top); we cite it rather than re-prove it.
Note that the value-preserving seal $\SVOpen(t) \to \SVClosed(t)$ does not
change any denotation under Definition~\ref{def:inst-sem} — (V2) reads the
bound value irrespective of open/closed — so among expansion's evolutions
only seal-to-top has semantic content, and it only \emph{adds} a binding
of value $\top$.

\subsection{The plumbing is sound}

\begin{lemma}[Worklist Containment]
\label{lem:plumbing}
Fix a mode $m$ and suppose the one-step containment
\[
\sem{i}_{\hat\rho} \subseteq \sem{E_m(i)}_{\hat\rho'}
\tag{$\ast$}
\]
holds for every instance $i$ popped during the run, where
$\hat\rho \to \hat\rho'$ is that step's evolution. Then
$\sem{t}_\rho \subseteq \sem{(I', C')}_{\rho'}$, and residual instances
appear in $I'$ verbatim.
\end{lemma}
\begin{proof}
Let $v \in \sem{t}_\rho = \sem{t.\mathsf{insts}}_\rho \cup
\sem{t.\mathsf{clauses}}_\rho$. All later evolutions are those of
Lemma~\ref{lem:evolutions}, so membership transports forward by
\textbf{L-MONO}; we may therefore evaluate all containments at the final
$\rho'$. If $v$ is justified by the clause part, it is in the initial
$C'$, and $\Clauses{::}\mathsf{unionN}$ only merges clauses keyed by
their var-sets using $\STypes$/clause union, which over-approximates set
union by \textbf{L-UNION} (paper B); the constructor invariant that a top
clause absorbs the rest, and \texttt{TType::create}'s dropping of $I'$
when $C'$ is top, only enlarge the result. If $v$ is justified by an
instance $i$, argue by induction on the number of worklist steps after
$i$ enters $\mathit{seen}$: when $i$ is popped, $(\ast)$ moves $v$ into
$\sem{E_m(i)}$, i.e.\ into its clauses (already covered) or one of its
instances $i''$; if $i''$ is fresh it enters the worklist and induction
applies; if $i'' \in \mathit{seen}$, then $i''$ was either already
processed (induction applies to its earlier occurrence — one-step
expansion of the same instance under the evolved state yields the same
$E_m(i'')$, by cache determinism for \texttt{TSInst}
(Theorem~\ref{thm:invariance}) and by value determinism of sealed reads
for the other arms) or is currently queued. The residual \texttt{TVar}
arm pushes $i$ itself onto $I'$; the top short-circuit returns $\top$
with $\sem{\top} = \Val \supseteq$ everything. \qedhere
\end{proof}

One genuine subtlety hides in the ``already processed'' case: a
\emph{recursive} named type re-emits its own \texttt{TSInst}, which the
$\mathit{seen}$ check then drops — the recursive occurrence contributes
nothing new only under a least-fixpoint reading of recursive definitions.
Section~\ref{sec:cycles} isolates this.

\begin{lemma}[mapSTypes Decomposition]
\label{lem:mapstypes}
\texttt{TType::mapSTypes}$(m, f)$ equals: expand in mode $m$; apply $f$
to each clause's \texttt{?STypes} (with a $\top$ expansion presented to
$f$ as the single top clause \texttt{None}); drop clauses whose image is
empty; rebuild with $\Clauses{::}\mathsf{unionN}$ and the untouched
residual $I'$. Consequently the main paper's refinement primitive
$\refop^{g,d}$ factors exactly as its Definition ``Memberwise
Refinement'' describes, with this paper supplying the expansion step and
paper B the rebuild unions; and \texttt{mapSTypes} with the identity
mapper is \texttt{expand}.
\impl{TType.sk / mapSTypes (388--414); TypeCheckRefine.sk /
setLocalFilterMap (600, 615)}
\end{lemma}
\begin{proof}
By inspection: \texttt{mapSTypes} calls \texttt{expandToVectors}, maps
$f$ through \texttt{Clause::create} (which returns \texttt{None} exactly
on an empty $\STypes$ image, implementing the drop), and rebuilds. The
top case at line 396--399 presents \texttt{None} to $f$ and produces
either the replacement clauses or $\top$. \texttt{rawMap} is the
substitution-free analogue used by generic transformers and performs no
expansion at all (its \texttt{fInsts} may convert instances to clauses
but has no access to a substitution); no refinement client reaches it,
so it carries no \textbf{L-EXP} obligation.
\impl{TType.sk / rawMap (362)}
\end{proof}

%==========================================================================
\section{Per-Construct Containment}
\label{sec:constructs}

\subsection{Type variables}
\label{sec:tvar}

\begin{lemma}[TVar One-Step Containment, modes $\NoneM$ and $\Cov$]
\label{lem:tvar}
For every variable occurrence $\alpha^K$ and state $\rho$:
\begin{itemize}
\item[(a)] Mode $\NoneM$: the step is the identity ($\alpha^K$ residual);
  containment and residual-semantics preservation are trivial.
\item[(b)] Mode $\Cov$, $\rho$ binds $\alpha$ (open or closed, freeze or
  not): the step returns $t \sqcap K$ built by \texttt{TType::inter} from
  the bound value $t$ and the occurrence constraints. By (V2), (V1) and
  \textbf{L-INTER} (paper B),
  $\sem{\alpha^K}_\rho \subseteq \sem{t}_\rho \cap \bigcap_K \sem{c}_\rho
  \subseteq \sem{t \sqcap K}_{\rho'}$, the last step reading the result
  at the post-step state via \textbf{L-MONO}. The seal
  $\SVOpen(t) \to \SVClosed(t)$, when it fires, is value-preserving.
\item[(c)] Mode $\Cov$, $\alpha$ unbound: the step seals
  $\alpha \mapsto \SVClosed(\top)$ and returns $\top \sqcap K$. By (V1)
  and \textbf{L-INTER},
  $\sem{\alpha^K}_\rho \subseteq \bigcap_K \sem{c}_\rho \subseteq
  \sem{\top \sqcap K}_{\rho'}$; the new binding of value $\top$ makes
  (V2) hold vacuously at $\rho'$.
\end{itemize}
\end{lemma}

\begin{finding}[Contravariant expansion does not satisfy L-EXP]
\label{find:contra}
We first attempted to prove the one-step containment for all three modes.
It fails in mode $\Contra$ on an \emph{unbound} variable:
$\getop(\alpha, \Contra)$ returns $\bot$ (no write), and
$\sem{\alpha^K}_\rho \subseteq \sem{\bot} = \emptyset$ is false whenever
the variable's envelope is inhabited. This is not a bug: a contravariant
read is a deliberate \emph{under}-approximation — ``the lower bound
gathered so far'' — consumed by the subtyping/growth machinery
(\texttt{growTVars}, \texttt{maybeGetRestElementType}), where the required
direction is reversed; that is the main paper's Obligation ``Validation
Soundness'', not this paper's. For a \emph{bound}
variable, case (b)'s argument goes through verbatim in mode $\Contra$,
since $\getop(\Contra)$ then returns the same bound value.
Consequence: axiom \textbf{L-EXP} must be scoped to modes
$\NoneM$/$\Cov$ (the mode scoping of Axiom \textbf{L-EXP};
Section~\ref{sec:status}).
\impl{TTypeSubst.sk / Subst::get (210--216); TTypeSubtyping.sk /
growTVars}
\end{finding}

\begin{lemma}[Refinement Clients Use Only Modes $\NoneM$/$\Cov$]
\label{lem:clients}
Every expansion performed by the refinement and constant-removal
machinery of the main paper uses mode $\Cov$ or mode $\NoneM$:
\begin{itemize}
\item \texttt{setLocalFilterMap} maps with \texttt{Some(Covariant())}
  (\texttt{TypeCheckRefine.sk:615});
\item the truthiness residues go through \texttt{residuePart} with
  \texttt{Some(Covariant())} (\texttt{TType.sk:350});
\item \texttt{conditionVerdict} and \texttt{getSTypes}-based readers use
  \texttt{expandToClauses} with \texttt{Covariant()}
  (\texttt{TType.sk:279, 631});
\item \texttt{removeConstant} expands with \texttt{None()}
  (\texttt{TType.sk:212}), leaving variables residual and untouched, as
  its model in the main paper's Definition ``Constant Removal'' requires
  (``the instance layer is unchanged'').
\end{itemize}
So the mode scoping of Finding~\ref{find:contra} costs the main paper
nothing.
\end{lemma}

\subsection{Named instantiations (\texttt{TSInst})}
\label{sec:tsinst}

$\instop(n, \bar{a})$ dispatches on the definition of $n$
(\texttt{TEnv::instantiate}):
a \texttt{TypeDef} binds each type parameter $P_i$ to the argument
$a_i$ \emph{as a type term} in the lexical \texttt{bindings} map
(checking constraints on a throwaway clone, diagnostics only) and
normalizes the alias body — a reference to $P_i$ in the body resolves to
the binding by \texttt{boundTy.withPos(pos)}, which changes only the
position and not the hash, so the argument term is inserted
\emph{verbatim}; a \texttt{Class} builds the object type via
\texttt{new}/\texttt{newObject}/\texttt{addFields} with the same
binding mechanism for field types; two early returns re-emit the
unexpanded $\mathsf{TSInst}$ itself — the synthetic missing-annotation
placeholder, and the recursion guard when $n$ is already being expanded
(\texttt{seenInstantiate}).
\impl{TEnv.sk / instantiate (110--148), instantiateTParams (174),
normalize TAInst arm (1354--1428), new (325), newObject (359)}

What normalization computes from the body is \texttt{TEnv} territory:
this paper does not model \texttt{normalize} and instead names the
interface assumption that the containment needs.

\begin{axiom}[E-NORM: Instantiation is Semantics-Preserving]
\label{ax:enorm}
For every named definition $n$ and arguments $\bar{a}$, and every state
$\rho$ at which the expansion runs,
$\sem{\mathsf{TSInst}(n,\bar{a})}_\rho \subseteq
 \sem{\instop(n,\bar{a})}_\rho$.
For the two early returns the result \emph{is} the input, so the axiom
is trivially true there and the re-emitted instance is a residual that
keeps its semantics by construction.
\impl{TEnv.sk / instantiate, normalize}
\end{axiom}

Under \textbf{E-NORM}, the \texttt{TSInst} one-step containment reduces
to one further claim: that reusing a \emph{cached} expansion — computed
at some earlier state $\rho_0$ — denotes at $\rho'$ no less than a fresh
$\instop$ run at $\rho'$ would. That is the centerpiece
(Section~\ref{sec:cache}); Theorem~\ref{thm:invariance} discharges it
for the alias path by showing the cached and fresh values are
\emph{equal}.

\subsection{Unexpanded intersections (\texttt{TInter})}
\label{sec:tinter}

The one-step expansion of $\mathsf{TInter}(\{t_1,\dots,t_k\})$ sets
\texttt{tenv.expandInter} and left-folds \texttt{TType::inter} over the
members. (That flag selects full over deferred field intersection; its
path-scoping and the deferred form's containment are registry row
TENV-EXPANDINTER (\textsc{open}), which this arm's argument imports.) By (I1) and \textbf{L-INTER} (paper B, which covers
\texttt{inter} at every level, including the \texttt{seenInter} cache
and the cycle guard \texttt{tenv.seenInter}),
\[
\sem{\mathsf{TInter}(\{t_i\})}_\rho = \bigcap_i \sem{t_i}_\rho
\subseteq \sem{\mathsf{interN}(\{t_i\})}_{\rho'},
\]
with the fold's intermediate evolutions transported by \textbf{L-MONO}.
When \texttt{inter\_} hits an in-progress pair it returns the
\emph{unexpanded} intersection \texttt{unexpandedInter}
(\texttt{TTypeInter.sk:154--156}); by (I1) that residual keeps its
semantics exactly. Similarly the bare-variable arms return a variable
occurrence with an added constraint — a residual whose semantics
(V1)-intersects the operand, again exact for our direction.
\impl{TTypeInter.sk / TInter::expand (37), inter (97), inter\_ (129),
unexpandedInter (82)}

\subsection{Delayed array literals (\texttt{TArrayLit})}
\label{sec:arraylit}

\texttt{ArrayLitTyping::expand} returns the committed type
$\mathsf{peek}(\alpha)$ — the binding of the literal's commitment
variable intersected with its constraints, read \emph{without} any state
transition — when $\alpha$ is bound, and otherwise the mutable tuple
type $\mathsf{Tuple}[\mathit{elts}]$ built from the widened element
types. The containment needs the commitment discipline documented at the
top of \texttt{TArrayLitTyping.sk} (``every commitment uses the target
type or the widened element types''), which is array-literal-typing
territory; we name it:

\begin{axiom}[E-ARR: Array-Literal Commitment Discipline]
\label{ax:earr}
For every $\mathsf{TArrayLit}(\alpha, \mathit{elts}, \_)$ and state
$\rho$: (i) if $\rho$ binds $\alpha$ then
$\sem{\mathsf{TArrayLit}(\ldots)}_\rho \subseteq
\sem{\mathsf{peek}(\alpha)}_\rho$ (the committed type contains the
literal's values); (ii) if $\alpha$ is unbound then
$\sem{\mathsf{TArrayLit}(\ldots)}_\rho \subseteq
\sem{\mathsf{Tuple}[\mathit{elts}]}_\rho$ (the literal's element values
inhabit the widened element types — a well-typedness fact about
how \texttt{TArrayLit} is constructed). Clause (ii) must hold at every
state, including states where $\alpha$ later becomes bound: the tuple
view taken by an early expansion must not be invalidated by a later,
wider commitment (e.g.\ to \texttt{Array<U>}).
\impl{TArrayLitTyping.sk / expand (100), arrayLitCommittedType (85)}
\impl{TArrayLitTyping.sk / arrayLitToTupleType (113); TTypeSubst.sk /
peek (224)}
\end{axiom}

\subsection{The containment theorem}

\begin{theorem}[T-EXP]
\label{thm:texp}
For $m \in \{\NoneM, \Cov\}$, all $t$, $\rho$: with
$(I', C', \rho') = \expandop_m(t, \rho)$,
\[
\sem{t}_\rho \subseteq \sem{(I', C')}_{\rho'},
\]
residual instances appear verbatim in $I'$ and keep their semantics,
and the state evolution $\rho \to \rho'$ consists only of the
transitions of Lemma~\ref{lem:evolutions}. The result is conditional on:
\textbf{L-UNION}/\textbf{L-INTER} (paper B), \textbf{L-MONO} (main
paper), \textbf{E-NORM}, \textbf{E-ARR}, and — for reuse of cached
\texttt{TSInst} expansions — Theorem~\ref{thm:invariance} on the alias
path and Obligation~\ref{obl:classsym} on the class path.
\end{theorem}
\begin{proof}
Lemma~\ref{lem:plumbing} reduces the claim to the one-step containment
$(\ast)$, established per construct: Lemma~\ref{lem:tvar} (variables),
Axiom~\ref{ax:enorm} plus Theorem~\ref{thm:invariance} (named
instantiations, cached or fresh), Section~\ref{sec:tinter}
(intersections), Axiom~\ref{ax:earr} (array literals).
\end{proof}

%==========================================================================
\section{Centerpiece: Soundness of the \texttt{seenExpansion} Cache}
\label{sec:cache}

The cache maps a $\mathsf{TSInst}$ to its expanded body and is consulted
before every $\instop$ call (Definition~\ref{def:expand}). The code
comment (\texttt{TTypeSubst.sk:17--31}) argues reuse stays sound as
inference proceeds because the cached body substitutes each type
parameter \emph{symbolically}: arguments appear as free variables,
re-resolved at every use through \texttt{Subst::get}; union/inter never
resolve an unbound \texttt{TVar}; and the one construct that reads a
binding during expansion, \texttt{keyof}, approximates identically
whether or not the variable is bound. This section makes those claims
precise, proves a strictly stronger statement for the alias path, and
corrects the comment in two places.

\subsection{Symbolic arguments, precisely}

\begin{definition}[Symbolic in the Arguments]
\label{def:symbolic}
$\instop(n, \bar{a})$ is \emph{symbolic in the arguments} if there is a
type-with-holes $B_n[\cdot_1,\dots,\cdot_k]$, determined by the
definition of $n$ alone, such that for every $\bar{a}$ and every state
$\rho$, $\instop(n,\bar{a})$ run at $\rho$ returns $B_n[\bar{a}]$ —
the same $\TType$ \emph{value} — where each hole occurrence is filled by
the argument term $a_i$ either (i) verbatim (up to
position, which is not hashed), or (ii) as a variable occurrence of
$a_i$'s variable with occurrence-level constraints added
(\texttt{setConstraint} — a functional update, registry row
SETCONSTRAINT-FUN), or (iii)
replaced by a \emph{binding-independent} approximation (a fixed type not
depending on $\rho$). Diagnostics emitted along the way may depend on
$\rho$; the returned value may not.
\end{definition}

\begin{lemma}[Symbolic $\Rightarrow$ Cache-Reuse Exactness]
\label{lem:symbolic-suffices}
If $\instop(n,\bar{a})$ is symbolic in the arguments, then for any
states $\rho_0$ (cache fill) and $\rho'$ (use), the cached value equals
what a fresh run at $\rho'$ would return; hence
$\sem{\text{cached}}_{\rho'} = \sem{\text{fresh}}_{\rho'}$ and the
\texttt{TSInst} one-step containment under reuse follows from
\textbf{E-NORM} at $\rho'$ alone.
\end{lemma}
\begin{proof}
Immediate: both runs return the value $B_n[\bar{a}]$. The point of the
definition is that argument \emph{bindings} are never consulted, so the
variables inside $\bar{a}$ remain free in the cached term and are
re-resolved through $\getop$ at every later expansion step of the
\emph{user} of the cached term — never frozen into the cache.
\end{proof}

\subsection{The alias path is binding-blind}

The mechanism turns out to be simpler and stronger than the comment's
operation-by-operation argument: on the \texttt{TypeDef} path,
normalization \emph{never consults the live substitution at all}.

\begin{definition}[Binding-Blind]
A computation is \emph{binding-blind} if every \texttt{Subst} whose
binding map it consults (via \texttt{get}, \texttt{peek},
\texttt{isUnbound}, \texttt{isOpen}, \texttt{isClosed}) is freshly
created with the empty state, or is a clone whose results influence
diagnostics only.
\end{definition}

\begin{lemma}[normalize is Binding-Blind]
\label{lem:blind}
\texttt{TEnv::normalize} and the helpers it reaches on the
\texttt{TypeDef} path are binding-blind. By arms of the dispatch
(\texttt{TEnv.sk:1104--1445}):
\begin{itemize}
\item \emph{Literals, primitives, top, bottom, tuples-as-names, and the
  fatal arms}: these construct constants or \texttt{TType::inst}; no
  substitution is involved.
\item \emph{\texttt{TAInst}} (1354): resolves through the lexical
  \texttt{bindings} map — a \texttt{TEnv} field, not a substitution —
  returning the bound term via \texttt{withPos} (verbatim insertion,
  case (i) of Definition~\ref{def:symbolic}); or emits an unexpanded
  \texttt{TType::inst}; the reference-site constraint check (1419--1426)
  runs on \texttt{Subst::mcreate(tcenvb)} — \emph{empty state} — and
  affects diagnostics only.
\item \emph{\texttt{TAUnion}/\texttt{TAInter}} (1429--1444): both create
  a fresh empty \texttt{Subst} (1430, 1438). \texttt{union} never reads
  bindings at any level (Lemma~\ref{lem:union-blind}). \texttt{inter}
  under the empty state sees every variable as unbound: bare-variable
  operands take the \texttt{setConstraint}/\texttt{unexpandedInter} arms
  (case (ii)); variables nested in compound operands are sealed to top
  \emph{in the throwaway state} and thus replaced by the
  binding-independent $\top$ (case (iii)) — see
  Finding~\ref{find:inter-comment}.
\item \emph{\texttt{TAKeyOf}} (1172--1175): creates a fresh empty
  \texttt{Subst} (1173) and expands the operand under it; see
  Finding~\ref{find:keyof}.
\item \emph{\texttt{TAObject}/\texttt{TAFun}/\texttt{TATypeOf}/tuples}:
  recursive normalization plus fresh empty substitutions at every
  operation site (\texttt{TEnv.sk:1451, 1556, 1567, 1582, 1639, 1752,
  1893}); the mapped-type key analysis (\texttt{FNIn}, 1953)
  expands the key type in mode $\NoneM$ under a fresh empty
  \texttt{Subst} (1893, 1973), leaving variables residual and unread.
\end{itemize}
\impl{TEnv.sk / normalize (1104), Subst creation sites (1173, 1421,
1430, 1438, 1451, 1556, 1567, 1582, 1639, 1752, 1893, 1973)}
\end{lemma}

\begin{theorem}[Alias-Path Expansion Invariance]
\label{thm:invariance}
For a \texttt{TypeDef} $n$, $\instop(n,\bar{a})$ is symbolic in the
arguments (Definition~\ref{def:symbolic}); in fact it returns the same
$\TType$ value at every state. Consequently (Lemma~\ref{lem:symbolic-suffices})
\texttt{seenExpansion} reuse is \emph{exact} on the alias path: the
cached body denotes, at every later state, exactly what a fresh
recomputation would.
\end{theorem}
\begin{proof}
$\instop$ on the \texttt{TypeDef} arm performs: parameter binding
(\texttt{instantiateTParams}) — the returned bindings are the argument
terms themselves, defaults normalized recursively, and the constraint
check runs on \texttt{subst.clone()} of the state \emph{copy} and is
discarded (\texttt{TEnv.sk:222}, diagnostics only); then
\texttt{normalize} of the body, which is binding-blind by
Lemma~\ref{lem:blind}. A binding-blind computation's result cannot
depend on the ambient binding map; the remaining inputs (definition
ASTs via \texttt{defMap}, the lexical bindings, the argument terms) are
fixed by $(n, \bar{a})$ and the compilation unit. The memo caches
consulted along the way (\texttt{seenInter}/\texttt{seenUnion}) belong
to the \emph{fresh empty} substitutions of Lemma~\ref{lem:blind} and
start empty at each site, so they cannot import state-dependence.
\end{proof}

\begin{corollary}[Immunity to the Freeze-Window Hazard]
\label{cor:freeze}
Papers B/E of this series identify a hazard family in the
\texttt{Subst} memo caches: under \texttt{freezeCandidates},
\texttt{Subst::get} returns an open candidate's bound \emph{without
sealing it}, so a cache entry whose value depends on that read
(\texttt{seenSub}, \texttt{seenInter}) can go stale when the candidate
later grows. On the alias path, \texttt{seenExpansion} is \emph{immune
by construction}: by Theorem~\ref{thm:invariance} the cached value does
not depend on any binding read, frozen or not — the unsealed read can
return whatever it likes without being able to reach the cached value.
(The class path is exactly where this immunity argument breaks;
Obligation~\ref{obl:classsym}.)
\end{corollary}

\begin{corollary}[Coherence Under save/restore/clone]
\label{cor:coherence}
\texttt{seenExpansion} is part of \texttt{SubstState}: \texttt{save},
\texttt{restore}, and \texttt{clone} carry bindings and cache
atomically, so an entry filled on an exploratory path is rolled back
with that path's bindings, and entries filled in a clone die with the
clone. On the alias path this atomicity is belt-and-braces — by
Theorem~\ref{thm:invariance} even a leaked entry would be correct — but
it is what keeps the \emph{class-path} exposure bounded to reuse within
one live substitution's lifetime.
\impl{TTypeSubst.sk / SubstState::restore (40), SavedSubst::create (63),
Subst::clone (137)}
\end{corollary}

\subsection{Verifying the comment's three claims}

\begin{lemma}[Union Never Resolves a TVar]
\label{lem:union-blind}
\texttt{TType::union} and everything it reaches perform no binding
reads: \texttt{union\_} unions the instance sets verbatim
(\texttt{TInsts::union}) and merges clauses by var-set key; the
recursive field/function cases (\texttt{TField::union},
\texttt{TFun::union}) recurse through \texttt{union} on covariant
components — and through \texttt{inter} on contravariant ones, which is
where the next lemma takes over. The \texttt{seenUnion} cache comment
(``any TVar in the operands passes through unchanged'') is confirmed at
the \texttt{TType} level.
\impl{TTypeUnion.sk / union (22), union\_ (42), TInsts::union (167)}
\end{lemma}

\begin{finding}[The comment's \texttt{inter} claim is imprecise]
\label{find:inter-comment}
We attempted to verify ``inter never resolves an unbound TVar''
verbatim and \emph{failed}: it is true only for \emph{bare} variable
operands (\texttt{inter\_}'s guard arms, \texttt{TTypeInter.sk:138--153},
which attach the other operand as an occurrence constraint). A variable
nested in a \emph{compound} operand — e.g.\
$(\alpha \mid \texttt{undefined})\ \&\ \texttt{string}$ — falls through
to \texttt{expandToClauses(Covariant())} on both operands
(\texttt{TTypeInter.sk:158--159}), which \emph{does} resolve $\alpha$:
sealing it to top if unbound, reading its value if bound. The claim is
rescued where the cache needs it by Lemma~\ref{lem:blind}: during alias
normalization \texttt{inter} runs under a fresh \emph{empty} state, so
the nested variable is always unbound there, the seal lands in the
throwaway state, and the substituted $\top$ is binding-independent —
and replacing a conjunct by $\top$ over-approximates the intersection,
so \textbf{E-NORM}'s direction is unharmed. Outside normalization (live
\texttt{inter} during checking) the compound-operand resolution is real
but is paper B's subject, not the cache's. Suggested action: reword the
comment's parenthetical to ``union caches independently of growth;
inter attaches \emph{bare} operand variables as constraints, and during
normalization runs on an empty substitution, so nested variables are
top-approximated, never bound-resolved.''
\end{finding}

\begin{finding}[The \texttt{keyof} claim holds vacuously — it never
reads]
\label{find:keyof}
The comment says \texttt{keyof} ``approximates an abstract or
unresolved operand identically (\texttt{string | number}) whether or not
the variable is bound.'' The audit shows something stronger: the
\texttt{TAKeyOf} arm creates a fresh \emph{empty} \texttt{Subst}
(\texttt{TEnv.sk:1173}) and \texttt{keyOfType} expands the operand under
it, so a variable argument is \emph{always} treated as unbound — its
covariant read seals it to top in the throwaway state, the expansion has
a top clause, and the guard at \texttt{TEnv.sk:2197--2206} returns
\texttt{string | number}. The live binding is unreachable, so cached and
fresh values agree not because the two reads approximate identically but
because there is no read. (The result also covers rigid parameters,
whose clauses have nonempty var-sets — the same guard.) Two caveats,
both outside the cache theorem: \texttt{keyOfType}'s own
over-approximation contract explicitly excludes validity checking
(``do not use keyof to check that a key is valid,''
\texttt{TEnv.sk:2188--2193}), and symbol keys of an unresolved operand
are not covered by \texttt{string | number} — both are
\textbf{E-NORM}-side concerns about what the approximation means, not
about cache invariance.
\impl{TEnv.sk / normalize TAKeyOf arm (1172--1175), keyOfType (2195)}
\end{finding}

\subsection{Where the proof stops: the class path}

\begin{obligation}[E-CLASS-SYM: Class Instantiation Symbolicity]
\label{obl:classsym}
On the \texttt{Class} arm, $\instop$ calls
\texttt{new} $\to$ \texttt{newObject} $\to$ \texttt{addFields} passing
the \emph{copy of the live state} made at \texttt{TEnv.sk:135} — not a
fresh empty substitution. Field-type normalization inside is
binding-blind (Lemma~\ref{lem:blind} applies arm by arm), but member
assembly is not obviously so: any \texttt{TType::inter} reached with
that substitution — the audited candidates are index merging
(\texttt{TField::interIndex}, \texttt{TTypeInter.sk:245}, called
unconditionally on \texttt{TType::inter}, unlike \texttt{TField::inter},
which defers via \texttt{unexpandedInter} when
\texttt{tenv.expandInter} is false), constructor assembly
(\texttt{findConstructor}), and \texttt{TIndexes::create} — can reach
\texttt{expandToClauses(Covariant())} on a compound operand mentioning a
type argument's variable and read its \emph{current} binding. Because
the read happens in the discarded copy, the seal that would normally
freeze the variable and justify the read (paper's
never-widen lemma (iii)) is thrown away: the live variable can still
grow, and the cached class body would then diverge from a fresh
recomputation — the same staleness family as the freeze-window hazard
(Corollary~\ref{cor:freeze}), reachable here even without
\texttt{freezeCandidates}, and additionally through it. \textbf{Status:
open.} We attempted both directions and completed neither: the proof
attempt dies at the \texttt{interIndex} chain above, and a concrete
end-to-end counterexample (a class instantiated at a bound-open
candidate whose member merge reads the bound, cached, followed by
observable growth) was not constructed. Discharge either by (a) showing
the reachable \texttt{inter} operands during \texttt{addFields} never
nest a live inference variable inside a compound type (targ variables
occur bare or inside unexpanded \texttt{TSInst}s there), or (b) making
the class path binding-blind by construction (fresh empty substitution
in \texttt{newObject}, as the alias path already effectively has), or
(c) keying \texttt{seenExpansion} entries for class instances by
substitution version, as \texttt{seenSub} already is for cycles.
\impl{TEnv.sk / instantiate (135, 141), newObject (359), addFields
(559)}
\impl{TTypeInter.sk / TField::interIndex (245), inter\_ (158)}
\end{obligation}

\begin{obligation}[E-GUARD: Recursion-Guard Identities Must Not Poison
the Live Cache]
\label{obl:guard}
$\instop$ returns the unexpanded $\mathsf{TSInst}$ itself when the
guard set \texttt{seenInstantiate} already contains $n$
(\texttt{TEnv.sk:128--130}), and \texttt{expandToVectors} caches
whatever $\instop$ returns (\texttt{TType.sk:653}). If such an
\emph{identity} entry were ever written to a long-lived substitution's
cache, every later expansion of that instance would hit the identity,
whose sole instance is then dropped by the $\mathit{seen}$ check — the
named type would contribute \emph{nothing} (effectively $\bot$) to every
future expansion, violating T-EXP outright. For the \emph{recursion-guard}
identity return this is empirically safe: \texttt{seenInstantiate} is
nonempty only inside an $\instop$ frame using the discarded substitutions
of Lemma~\ref{lem:blind}, and a recursive alias's inner \texttt{TSInst}
sits nested in a field/argument the head worklist never processes
(\texttt{type A = A | number} reads as \texttt{number}). \textbf{Status
of the recursion-guard case: open (small)}. But there is a \emph{second}
identity return — the missing-annotation placeholder
(\texttt{TEnv.sk:122--124}) — that fires with \texttt{seenInstantiate}
\emph{empty}, at a top-level live-cache-fill site; the re-emitted identity
equals the popped instance and is dropped by the expansion seen-set,
leaving empty clauses, so an un-annotated ambient-class field expands to
$\bot$ and reads into any type including \texttt{never}. That is
registry row \texttt{E-GUARD-PLACEHOLDER}, \textsc{violated} and
witnessed end-to-end (issue draft
\texttt{ambient-unannotated-field-bottom.md}); the fix keeps an identity
\texttt{instantiate} result as a residual (or returns a top clause)
rather than letting the seen-set drop it. A cache-write invariant that only asserts
\texttt{seenInstantiate} is empty does \emph{not} catch the placeholder
(it is empty there); the invariant must instead assert ``the cached
value's insts $\neq$ the key.''
\impl{TEnv.sk / instantiate (125--130); TType.sk / expandToVectors
(646--656)}
\end{obligation}

%==========================================================================
\section{Alias Cycles and Termination}
\label{sec:cycles}

Three guards keep expansion of recursive definitions finite; soundness of
each cutoff is the point here, termination secondary.

\paragraph{One-level unfolding (\texttt{seenInstantiate}).}
While expanding $n$, an inner reference to $n$ normalizes to an
unexpanded \texttt{TSInst} instead of unfolding
(\texttt{TEnv.sk:45--60, 128--131, 1405--1407}). The cutoff is sound by
construction: the re-emitted instance denotes the named type ((N1)),
and nested occurrences sit inside object fields or targs, where the
worklist does not descend — expansion is one level deep by design (the
main paper's ``expansion rewrites a type to its instance-plus-clauses
view'' is inherently head expansion).

\paragraph{The worklist $\mathit{seen}$ set.}
A \emph{top-level} recursive occurrence — \texttt{type A = A | number}
normalizes to $(\{\mathsf{TSInst}(A)\}, \{\texttt{number}\})$ — is
deduplicated: the body's $\mathsf{TSInst}(A)$ equals the instance being
expanded and is dropped (\texttt{TType.sk:671--675}). The expansion of
$A$ is then just \texttt{number}. This is sound under a least-fixpoint
reading of type definitions: the equation
$X = X \cup \sem{\texttt{number}}$ has least solution
$\sem{\texttt{number}}$, and any value's membership in a recursively
defined type must be justified by finite unfolding. We record this
reading as part of (N1)'s meaning rather than as a separate assumption;
a definition intended coinductively (values justified only by infinite
regress) has no inhabitants a program can construct.
Mutual cycles behave the same: expanding \texttt{type A = B; type B = A}
yields empty clauses (bottom) after both identities deduplicate — the
least fixpoint — while the checker independently rejects such
definitions at validation (TS2456).

\paragraph{Fatal circularity (\texttt{forbiddenInsts}).}
Within one $\instop$ frame, a \emph{second} reference to the
in-expansion name is fatal (TS2456, \texttt{TEnv.sk:1375--1379}): the
first reference marks the name forbidden (1405--1407), the next lookup
dies. Generative loops (\texttt{type A<T> = A<A<T>> | number}) die the
same way — normalizing the inner argument \texttt{A<T>} is a second
reference. A fatal diagnostic ends checking, which is trivially sound.
Termination of the worklist itself: each iteration either pops a
$\mathit{seen}$ instance's expansion products (each distinct instance is
processed at most once) or returns $\top$; non-termination would require
$\instop$ to keep manufacturing distinct instances at the top level of
bodies, which the generative-loop fatal cuts off on the alias path. We
have not proved a global termination bound for the class path (chains of
distinct \texttt{TSInst}s through superclass clauses); termination is
outside the soundness contract and is left unstated.
\impl{TEnv.sk / seenInstantiate (60), forbiddenInsts (64), normalize
TAInst (1375, 1405); TType.sk / expandToVectors (671--675)}

%==========================================================================
\section{Status and Ledger Edits}
\label{sec:status}

\subsection{Decomposition of L-EXP}

\begingroup
\renewcommand{\arraystretch}{1.25}
\small
\begin{longtable}{@{}p{2.4cm}p{7.8cm}p{3.3cm}@{}}
\toprule
\textbf{Item} & \textbf{Statement} & \textbf{Anchor} \\
\midrule
\endfirsthead
\toprule
\textbf{Item} & \textbf{Statement} & \textbf{Anchor} \\
\midrule
\endhead
\bottomrule
\endlastfoot
E-PLUMB & Worklist, clause rebuild, \texttt{mapSTypes} factoring, top
  short-circuit, residual pass-through
  (Lemmas~\ref{lem:plumbing}, \ref{lem:mapstypes}) &
  \texttt{expandToVectors}, \texttt{mapSTypes} \\
E-TVAR & Variable one-step containment, modes $\NoneM$/$\Cov$
  (Lemma~\ref{lem:tvar}); mode $\Contra$ excluded with counterexample
  (Finding~\ref{find:contra}); refinement clients verified
  $\NoneM$/$\Cov$-only (Lemma~\ref{lem:clients}) &
  \texttt{Subst::get} \\
E-INTER & \texttt{TInter} one-step containment; cycle residual exact
  (\S\ref{sec:tinter}) & \texttt{TInter::expand} \\
E-NORM & Instantiation is semantics-preserving at the state of use
  (Axiom~\ref{ax:enorm}) — the \texttt{TEnv}/normalize interface &
  \texttt{TEnv::instantiate} \\
E-ARR & Array-literal commitment discipline (Axiom~\ref{ax:earr}) &
  \texttt{TArrayLitTyping.sk} \\
E-SYM-ALIAS & Cache invariance, \texttt{TypeDef} path: binding-blind,
  hence exact reuse and freeze-window immunity
  (Theorem~\ref{thm:invariance}, Corollary~\ref{cor:freeze}) &
  \texttt{seenExpansion}; \texttt{TEnv::normalize} \\
E-CLASS-SYM & Cache invariance, \texttt{Class} path
  (Obligation~\ref{obl:classsym}) & \texttt{TEnv::newObject};
  \texttt{TField::interIndex} \\
E-GUARD & Recursion-guard identities never cached into a long-lived
  substitution (Obligation~\ref{obl:guard}) &
  \texttt{expandToVectors} (653) \\
\end{longtable}
\endgroup

Net effect on \textbf{L-EXP}: \textsc{assumed} $\to$ \textsc{proved}
for modes $\NoneM$/$\Cov$, conditional on
\textbf{L-UNION}/\textbf{L-INTER}, \textbf{L-MONO}, \textbf{E-NORM},
\textbf{E-ARR}, \textbf{E-CLASS-SYM}, \textbf{E-GUARD}. The unconditional
residue is smaller and sharper than the original axiom: two named
interface assumptions and two localized open obligations.

\subsection{Registry}

The authoritative, current statuses are maintained in the registry
(\texttt{docs/ledger.yaml}), rendered as the main paper's Assumptions and
Proof Obligations tables; this paper is cited there as evidence — in
particular for the mode scoping of Axiom \textbf{L-EXP} (the containment
holds for modes \texttt{None}/\texttt{Covariant} only,
Lemma~\ref{lem:clients}) and for the rows E-NORM, E-ARR, E-CLASS-SYM,
E-GUARD. The registry rows this paper supports:

% GENERATED by docs/tools/docs.py from docs/ledger.yaml — do not edit.

\begin{center}
\renewcommand{\arraystretch}{1.25}
\small
\begin{tabular}{@{}>{\raggedright\arraybackslash}p{2.4cm}>{\raggedright\arraybackslash}p{8.2cm}>{\raggedright\arraybackslash}p{2.6cm}@{}}
\toprule
\textbf{Row} & \textbf{Statement} & \textbf{Status} \\
\midrule
L-EXP & Expansion preserves or enlarges semantics (modes \texttt{None}/Covariant); residual instances keep theirs. & \textsc{proved} \\
E-NORM & Named-type normalization is semantics-preserving (interface consumed by expansion). & \textsc{assumed} \\
E-ARR & Delayed array-literal forms keep their semantics under expansion. & \textsc{assumed} \\
E-CLASS-SYM & Class expansion either does not depend on existing variable bindings, or any bindings it reads cannot make cache reuse unsound. & \textsc{open} \\
E-GUARD & Recursion-guard identity expansions are never cached into long-lived substitutions. & \textsc{open} \\
E-PLACEHOLDER-ATOM & An un-annotated ambient-class field reads as an opaque per-field atom; the atom's denotation must over-approximate the arbitrary runtime value. & \textsc{open} \\
E-EXPKEY & The expansion-cache key (\texttt{ExpansionKey}: unqualified name text + arguments) identifies one definition: same-named, same-argument aliases from different modules must not collide. & \textsc{violated} \\
\bottomrule
\end{tabular}
\end{center}


\paragraph{Imports (review contract).}
The named statements this paper assumes from elsewhere. A reviewer treats
each as an axiom, checking citations against the registry's canonical
statement; no other paper's internals are in scope:

% GENERATED by docs/tools/docs.py from docs/ledger.yaml — do not edit.

\begin{center}
\renewcommand{\arraystretch}{1.25}
\small
\begin{tabular}{@{}>{\raggedright\arraybackslash}p{2.6cm}>{\raggedright\arraybackslash}p{7.4cm}>{\raggedright\arraybackslash}p{1.2cm}>{\raggedright\arraybackslash}p{2.2cm}@{}}
\toprule
\textbf{Row} & \textbf{Canonical statement} & \textbf{Owner} & \textbf{Status} \\
\midrule
L-UNION & \texttt{union} over-approximates set union (all levels). & B & \textsc{proved} \\
L-INTER & \texttt{inter} over-approximates set intersection (the containment clause). & B & \textsc{proved} \\
L-MONO & Refinement-time state evolutions preserve membership in already-installed types. & E & \textsc{proved} \\
TVAR-SEM & Sharpened variable denotation (get-time constraint intersection as definition), verified against every consumer. & B & \textsc{open} \\
SETCONSTRAINT-FUN & \texttt{TVar::setConstraint} is a functional update on a frozen value: constraints are occurrence-level, not shared state. & E & \textsc{proved} \\
TENV-EXPANDINTER & The \texttt{tenv.expandInter} mode flag is path-scoped, and both intersection forms it selects (expanded and deferred \texttt{unexpandedInter}) satisfy the L-INTER containment. & E & \textsc{open} \\
\bottomrule
\end{tabular}
\end{center}


\paragraph{Outstanding recommendations.}
(i) Fix the code comment \texttt{TTypeSubst.sk:17--31} per
Finding~\ref{find:inter-comment} (inter's ``never resolve'' holds for
bare operands; nested variables are top-approximated because
normalization substitutions are empty) and Finding~\ref{find:keyof}
(keyof never reads the live binding — \texttt{TEnv.sk:1106}).
(ii) Wire the E-GUARD pinned test; record
Lemma~\ref{lem:evolutions}'s inventory (expansion performs
get-transitions and cache writes only; never a grow) as the concrete
content behind L-MONO's ``evolutions of the kind covered''.

\subsection{Findings from failed proof attempts (index)}

Finding~\ref{find:contra}: unconditional (all-modes) L-EXP is false;
counterexample unbound-variable contravariant read returning $\bot$.
Finding~\ref{find:inter-comment}: the comment's blanket
``inter never resolves an unbound TVar'' failed verification; rescued
only via the binding-blind property.
Finding~\ref{find:keyof}: the keyof-equal-approximation claim could not
be proved as stated (a bound variable read \emph{would} produce precise
keys) — it is true for the vacuous reason that no live read occurs.
Obligation~\ref{obl:classsym}: class-path symbolicity resisted both
proof (dies at the \texttt{interIndex} chain) and refutation (no
end-to-end counterexample constructed).
Obligation~\ref{obl:guard}: cache poisoning by recursion-guard
identities could not be excluded by a local invariant; global call-site
audit only.

%==========================================================================
\section*{Implementation Anchors (Summary)}

\begin{center}
\renewcommand{\arraystretch}{1.2}
\footnotesize
\begin{tabular}{@{}l@{\quad}l@{}}
\toprule
\textbf{Model} & \textbf{Implementation (rev.\ \texttt{\skjspinhash})} \\
\midrule
$\expandop_m$, worklist & \texttt{TType.sk / expandToVectors (634)} \\
rebuilds & \texttt{expand (680), expandToClauses (692), mapSTypes (388),
  rawMap (362)} \\
$\getop$ transitions, freeze & \texttt{TTypeSubst.sk / Subst::get (188),
  freezeCandidates (89)} \\
cache + saved state & \texttt{TTypeSubst.sk / seenExpansion (31),
  restore (40), clone (137)} \\
cache read/write site & \texttt{TType.sk / expandToVectors (646--656)} \\
$\instop$, parameter binding & \texttt{TEnv.sk / instantiate (110),
  instantiateTParams (174)} \\
symbolic resolution & \texttt{TEnv.sk / normalize TAInst arm
  (1354--1373), bindings (43)} \\
fresh empty substitutions & \texttt{TEnv.sk / 1173, 1421, 1430, 1438,
  1451,} \\
 & \texttt{1556, 1567, 1582, 1639, 1752, 1893, 1973} \\
keyof & \texttt{TEnv.sk / keyOfType (2195)} \\
mapped types & \texttt{TEnv.sk / FNIn arm (1953)} \\
class path & \texttt{TEnv.sk / new (325), newObject (359), addFields
  (559)} \\
intersections & \texttt{TTypeInter.sk / TInter (17), expand (37),
  inter\_ (129), interIndex (245)} \\
unions & \texttt{TTypeUnion.sk / union (22), union\_ (42)} \\
array literals & \texttt{TArrayLitTyping.sk / expand (100)} \\
recursion guards & \texttt{TEnv.sk / seenInstantiate (60),
  forbiddenInsts (64)} \\
\bottomrule
\end{tabular}
\end{center}

\end{document}
